\section{LIT-AB: John List Research: Field Experiments and Market Behavior}
\label{app:litab}

This appendix integrates papers by John List on field experiments, market behavior, and allocation
mechanisms. List's research program demonstrates that laboratory behavioral anomalies often vanish
in field settings with real stakes and experienced participants, fundamentally challenging behavioral
economics claims about market efficiency.

\subsection{Research Program Overview}

John List's research addresses core behavioral economics and field experimental themes:

\begin{itemize}
  \item \textbf{Field Experiments}: Testing economic behavior in naturally occurring settings
  \item \textbf{Market Efficiency}: When and why markets eliminate behavioral anomalies
  \item \textbf{Experience Effects}: How experience and stakes affect decision-making
  \item \textbf{Auction Design}: How auction mechanisms affect bidding and efficiency
  \item \textbf{Environmental Valuation}: Field methods for non-market good valuation
  \item \textbf{Charitable Behavior}: Mechanisms to increase giving and pro-social behavior
  \item \textbf{Mechanism Design}: Using field experiments to test institutional design
  \item \textbf{Policy Applications}: Using field evidence for evidence-based policy
\end{itemize}

\subsection{Papers Integrated: 0 works}


\subsection{Summary}

This appendix integrates 0 papers by John List spanning research on
field experiments, market efficiency, environmental valuation, auction design, and policy applications.
List's research demonstrates the power of field experiments to test economic theories in real-world settings
and design effective policies.

Key findings from List's research:
\begin{enumerate}
  \item Behavioral anomalies observed in labs often disappear in field settings
  \item Market experience and real stakes eliminate many behavioral effects
  \item Field experiments provide powerful evidence for policy design
  \item Mechanism design works better when tested in the field
  \item Environmental valuation methods show large disparities
  \item Simple interventions (asking, matching) can significantly increase pro-social behavior
\end{enumerate}

\subsection{References}

For complete reference details and citations, see the master bibliography
\texttt{bcm2\_master\_references.bib}.

\vspace{1em}
\hrule
\vspace{0.5em}

\noindent\textit{Cross-references:}
Appendix G (Central Glossary), Chapter 14 (Applications)
